\section{Introduction}
\label{sec:intro}

% --- Motivation ---
% 1. Test-time scaling is the dominant recipe for "make the LLM reason better".
% 2. MCTS is the most elaborate (and most expensive) member of that family.
% 3. Most evidence is confounded: more search == more compute.

Test-time search has become a default recipe for coaxing more reliable reasoning
out of large language models (LLMs): sample several chains and vote
(self-consistency), score samples and keep the best (best-of-$N$), expand a beam,
or run Monte Carlo Tree Search (MCTS) over intermediate steps. Reported gains,
however, are frequently confounded: a search method issuing dozens of model calls
is compared against a single greedy call, so it is unclear whether the
improvement comes from \emph{search structure} or simply from \emph{spending more
compute}.

\begin{figure}[t]
\begin{lstlisting}
// gold: CWE-416 (Use After Free)
   int new_high = tree[current].high;
   tree[current].high = low-1;
   add_range(ctx, cmap, high+1, new_high, ...);
+  tree = cmap->tree;   // fix: re-fetch; add_range may realloc
 }
 if (tree[current].low > high) {
\end{lstlisting}
\caption{A motivating pilot example (ghostscript; abridged). \texttt{add\_range}
may reallocate \texttt{cmap->tree}, so the subsequent reads through the stale
\texttt{tree} pointer are a use-after-free (CWE-416). Greedy answers CWE-125
(out-of-bounds \emph{read}); self-consistency at $8\times$ the budget agrees with
greedy; MCTS over the taxonomy spends $10.7$k tokens, commits at the first level
to the \emph{spatial} out-of-bounds category, and lands on CWE-787 --- the wrong
branch, unrecoverable. MCTS over a \emph{reasoning trajectory} notices the
realloc-then-reuse pattern and answers CWE-416 at a comparable budget. More
compute did not help; searching a different space did.}
\label{fig:motivating}
\end{figure}

% --- Problem framing ---
We study this question on \emph{hierarchical vulnerability classification}: given
a code fragment, predict a root-to-leaf path in the CWE taxonomy. This task is a
good testbed because (i) it has a real, structured output space with a natural
search tree; (ii) errors compound across levels, so search has a plausible
mechanism to help \emph{or} to hurt; and (iii) public, CWE-labeled datasets exist.

% --- The confound we remove ---
\paragraph{The compute confound.}
Our central methodological commitment is \textbf{compute-matched comparison}: every
strategy is evaluated across a grid of token budgets $\budget$, and each method's
knob ($N$, beam width, or MCTS iterations) is tuned so its expected token spend
matches $\budget$. The headline result is an accuracy-vs-budget curve, not a single
``MCTS beats greedy'' number.

% --- Two axes the literature entangles ---
\paragraph{Three axes we isolate.}
Whether and how search helps plausibly depends on (a) the base policy's strength,
(b) the fidelity of the reward guiding it, and (c)---we argue, most
importantly---the \emph{space} the search operates in. We vary \emph{policy
strength} $\polacc$ across LLMs, contrast self-evaluation against a stronger
external judge, and compare searching the output \emph{taxonomy} against searching
a free-form \emph{reasoning trajectory}.

% --- Contributions ---
\paragraph{Contributions.}
\begin{enumerate}
\item Our central result: \textbf{what you search over dominates how much you
search}. MCTS over a reasoning trajectory beats MCTS over the taxonomy by $+12$
leaf-accuracy points at \emph{half} the token budget ($p{<}10^{-4}$), and its
accuracy is flat in the iteration knob---the reasoning \emph{scaffold}, not the
search machinery, carries the gain. A cascade analysis shows why: taxonomy search
commits to a wrong coarse category on $60\%$ of examples and cannot recover,
whereas reasoning search defers commitment (\S\ref{sec:results}).
\item A \textbf{compute-matched protocol} for evaluating test-time search on a
structured-output task, removing the search-vs-compute confound, with a controlled
comparison of strategies (greedy, self-consistency, best-of-$N$, beam, $\mctstax$,
$\mctsreas$) on public CWE-labeled data (\S\ref{sec:experiments}).
\item The marginal gain of search \emph{decreases} with policy strength: large for
weak policies, negligible-to-negative for strong ones---test-time search
compensates for weak policies rather than improving strong ones
(\S\ref{sec:results}).
\item A stronger external-judge reward does not improve search, locating the
bottleneck in the policy and the search structure rather than reward fidelity
(\S\ref{sec:results}).
\end{enumerate}

% --- Honest framing ---
We emphasize that this is a \emph{diagnostic} study: a finding that MCTS does not
beat cheaper search at matched budget is a valid and useful result, and we design
the paper so its contribution does not depend on the sign of that result.
